% Lazo de control de la dosificadora de harina de Rumapea et al. (2026).
% Elaboración propia. Línea continua: harina. Línea discontinua: señales.
% Flujo fino por reducción del ángulo de la compuerta; cierre a 0° en la consigna (p. 48).
\begin{tikzpicture}[
    >=Latex,
    font=\small,
    bloque/.style={
        rectangle,
        draw,
        rounded corners=2pt,
        minimum height=0.9cm,
        text width=3.4cm,
        align=center,
        inner sep=3pt
    },
    material/.style={
        -{Latex[width=6pt,length=5pt]},
        thick
    },
    senal/.style={
        -{Latex[width=6pt,length=5pt]},
        thick,
        dashed
    }
]

% Recorrido de la harina (x = 0)
\node[bloque] (tolva) at (0, 3.0) {Tolva cónica};
\node[bloque] (comp)  at (0, 1.5) {Compuerta con servomotor};
\node[bloque] (plat)  at (0, 0.0) {Plataforma sobre celda de carga de 1~kg};

% Control (x = 7)
\node[bloque] (cons) at (7.0, 3.0) {Consigna $m_{\mathrm{c}}$};
\node[bloque, text width=6.2cm, minimum height=2.6cm] (ctrl) at (7.0, 0.75) {%
    \textbf{Controlador (Arduino Uno)}\par\vspace{3pt}
    $m < 0{,}9\,m_{\mathrm{c}}$: apertura completa\\
    $0{,}9\,m_{\mathrm{c}} \le m < m_{\mathrm{c}}$: ángulo reducido\\
    $m \ge m_{\mathrm{c}}$: cierre a 0°};

% Material (continua) y señales (discontinua)
\draw[material] (tolva) -- (comp);
\draw[material] (comp) -- (plat);
\draw[senal] (plat.east) -- node[below, font=\footnotesize] {masa $m$}
    (plat.east -| ctrl.west);
\draw[senal] (comp.east -| ctrl.west) -- node[above, font=\footnotesize]
    {apertura} (comp.east);
\draw[senal] (cons) -- (ctrl);

\end{tikzpicture}
